\chapter[可解释性、归因公理与算法公平性]{机械可解释性、归因公理\\与算法公平性伦理}\label{ch:interp-fairness}

\chapEpigraph{你的行动，要把你自己人身中的人性，和其他任何人身中的人性，在任何时候都同样看作是目的，永远不能只看作是手段。}{Immanuel Kant（伊曼努尔·康德，《道德形而上学奠基》）}

\begin{chapIntro}{本章核心问题与学习地图}
当一个深度神经网络在测试集上取得 $95\%$ 的准确率时，我们距离把它部署到临床诊疗、信贷审批或司法风险评估还有多远？答案是：非常遥远。高准确率只回答了“它常不常猜对”，却没有回答两个决定系统能否赢得社会信任的根本问题：\textbf{第一，它究竟凭借什么依据做出这个判断？第二，它的错误是否以系统性的不公方式落在特定人群身上？}

本章从“黑箱归因”到“白箱机理”，再到“数学伦理”，构建可解释性与算法公平性的严谨理论框架：
\begin{enumerate}
  \item \textbf{梯度归因的饱和失效与公理化诉求}：用显式反例证明原始梯度 $\nabla_x F(x)$ 违背敏感性公理，明确特征归因必须遵守的公理标尺；
  \item \textbf{积分梯度（IG）与完备性公理证明}：通过多元链式法则与微积分基本定理严格证明 $\sum_i \IG_i(x; x') = F(x) - F(x')$，推导中点黎曼和的 $O(1/m^2)$ 收敛阶并在实测 CSV 上验证；
  \item \textbf{合作博弈论与 SHAP / Shapley 值}：给出精确 Shapley 唯一性公式，通过三特征手算实例核验效率、对称、哑元与可加四公理，并剖析 KernelSHAP 与 TreeSHAP；
  \item \textbf{机械可解释性（Mechanistic Interpretability）}：辨析线性探针的记忆幻觉，建立激活修补（Activation Patching）因果追踪范式，以及利用稀疏自编码器（SAE）解耦神经叠加（Superposition）；
  \item \textbf{算法公平性与 Chouldechova 不可能定理}：定义统计均等、均等赔率与预测均等，用贝叶斯定理严格证明基础率不等时“校准”与“误报/漏报均等”不可兼得的代数恒等式，并以实测数据验证。
\end{enumerate}
\end{chapIntro}

\section{为什么原始梯度归因会失效？敏感性与实现不变性公理}

设神经网络的标量预测函数为 $F: \R^d \to \R$（例如某目标类别的 logit 或概率），输入特征向量为 $x = (x_1, \dots, x_d)^\top \in \R^d$。要解释模型在点 $x$ 处的预测，人们首先需要一个\textbf{参照基线（Baseline）} $x' \in \R^d$（代表“缺少信息”的中性输入，如全黑图像、全零词嵌入或特征均值）。所谓\textbf{特征归因方法（Attribution Method）}，就是给出向量 $A_F(x, x') = (a_1, \dots, a_d)^\top \in \R^d$，其中 $a_i$ 量化第 $i$ 个特征对输出增量 $F(x) - F(x')$ 的贡献。

{\rightskip=0pt plus 1em\relax 最直观的候选方法莫过于计算局部梯度 $\nabla_x F(x)$，或将其乘以输入差值 $(x_i - x'_i)\,\frac{\partial F(x)}{\partial x_i}$（即 \texttt{Input} $\times$ \texttt{Gradient}）。然而，Sundararajan 等人（ICML 2017）指出，任何合格的归因方法至少应当满足两条最基本的公理，而原始梯度在第一条上便彻底翻车。\par}

\begin{defBox}
\begin{definition}[归因的敏感性公理与实现不变性公理]\label{def:ig-axioms}
设 $F: \R^d \to \R$ 为预测函数，$x, x' \in \R^d$ 分别为待解释输入与基线：
\begin{enumerate}
  \item \textbf{敏感性（a）（Sensitivity(a)）}：若输入 $x$ 与基线 $x'$ 仅在单一特征 $i$ 上取值不同（$x_i \neq x'_i$，而 $\forall j \neq i,\; x_j = x'_j$），且函数输出确实发生了变化（$F(x) \neq F(x')$），则该特征的归因值必须非零：$a_i \neq 0$。
  \item \textbf{实现不变性（Implementation Invariance）}：若两个网络 $F_1, F_2$ 对所有输入都定义完全相同的数学函数（$\forall x \in \R^d,\; F_1(x) = F_2(x)$），哪怕它们内部网络结构、隐藏层数或参数化方式截然不同，两者的归因结果必须完全一致：$A_{F_1}(x, x') = A_{F_2}(x, x')$。
\end{enumerate}
\end{definition}
\end{defBox}

\begin{example}[单变量饱和反例：原始梯度为何对真正起作用的特征“视而不见”？]\label{ex:grad-fail}
考虑极其简单的一维 ReLU 网络：
\[
  F(x) = 1 - \operatorname{ReLU}(1 - x) =
  \begin{cases}
    x, & x \le 1, \\
    1, & x > 1.
  \end{cases}
\]
取基线 $x' = 0$，输入 $x = 2$。显然输出从 $F(0) = 0$ 剧烈跃升到 $F(2) = 1$，且整个跃升完全归功于唯一输入特征 $x$ 从 $0$ 变到了 $2$。然而，在端点 $x = 2$ 处函数已经进入\textbf{饱和平坦区（Saturation Regime）}：
\[
  \left.\frac{\partial F}{\partial x}\right|_{x=2} = 0
  \quad\Longrightarrow\quad
  (x - x')\left.\frac{\partial F}{\partial x}\right|_{x=2} = (2 - 0)\cdot 0 = 0!
\]
原始梯度与 \texttt{Input} $\times$ \texttt{Gradient} 竟然将归因值判为 $0$，公然违背敏感性公理！在真实深度网络中，Sigmoid、Tanh、Softmax 与带截断的 ReLU 组合极易在高置信度样本附近形成局部平坦饱和区，导致单点梯度“装瞎”，把关键决策特征误判为无关特征。
\end{example}

同样值得警惕的是，早期视觉解释界流行的 DeConvNet、Guided Backpropagation 与部分离散传播规则（如 LRP、DeepLIFT）在反向传播时人为修改了非线性门的传播规则，不再只依赖函数 $F$ 的输入--输出映射，而是依赖于网络的具体层拆解方式，因而违背了\textbf{实现不变性}：只需对网络做一次不改变数学函数的等价权重重构，解释热力图就会面目全非。

\section{积分梯度（IG）：完备性公理证明与数值收敛实测}

为什么局部梯度在例 \ref{ex:grad-fail} 中失败？因为它只偷看了终点 $x = 2$ 一瞬间的斜率，却忽略了输入从基线 $x' = 0$ 一路走到 $x = 2$ 的全过程中，$x \in [0, 1]$ 区间内斜率始终等于 $1$ 的事实！如果我们将从基线 $x'$ 沿直线移动到输入 $x$ 的整条路径上的所有梯度累积起来，便得到了优雅的\textbf{积分梯度（Integrated Gradients, IG）}。

\begin{defBox}
\begin{definition}[积分梯度，Sundararajan et al., ICML 2017]\label{def:ig}
设函数 $F: \R^d \to \R$ 在连接基线 $x'$ 与输入 $x$ 的直线路径
\[
  \gamma(\alpha) = x' + \alpha (x - x'), \qquad \alpha \in [0, 1]
\]
的某个开邻域内几乎处处可微（且沿路径绝对连续）。则第 $i$ 个特征的\textbf{积分梯度}定义为
\begin{equation}\label{eq:ig-def}
  \IG_i(x; x') \;=\; (x_i - x'_i) \int_0^1 \frac{\partial F\big(x' + \alpha(x - x')\big)}{\partial x_i}\,d\alpha,
  \qquad i = 1, \dots, d.
\end{equation}
\end{definition}
\end{defBox}

\begin{recallBox}{多元链式法则与微积分基本定理（第 1 卷第 4 章）}
若标量函数 $F: \R^d \to \R$ 与曲线 $\gamma: [0,1] \to \R^d$ 均可微（或分段光滑且绝对连续），则复合单变量函数 $g(\alpha) = F(\gamma(\alpha))$ 的导数为方向导数：
\[
  g'(\alpha) = \inner{\nabla F(\gamma(\alpha))}{\gamma'(\alpha)} = \sum_{i=1}^d \frac{\partial F(\gamma(\alpha))}{\partial x_i}\,\gamma'_i(\alpha),
\]
且满足微积分基本定理 $\int_0^1 g'(\alpha)\,d\alpha = g(1) - g(0)$。
\end{recallBox}

\begin{thmBox}
\begin{theorem}[积分梯度的完备性公理与导出性质]\label{thm:ig-completeness}
在定义 \ref{def:ig} 的条件下，积分梯度严格满足\textbf{完备性公理（Completeness Axiom）}：所有特征的归因之和精确等于预测输出相对于基线的总增量，即
\begin{equation}\label{eq:ig-completeness}
  \sum_{i=1}^d \IG_i(x; x') \;=\; F(x) - F(x').
\end{equation}
进而，积分梯度自动满足：
\begin{enumerate}
  \item \textbf{敏感性（a）}：若仅特征 $i$ 变化且 $F(x) \neq F(x')$，则其余项皆因 $x_j - x'_j = 0$ 而为零，必有 $\IG_i(x; x') = F(x) - F(x') \neq 0$；
  \item \textbf{实现不变性}：定义式仅涉及函数 $F$ 的偏导数，与内部网络实现无关；
  \item \textbf{线性性}：对任意 $a, b \in \R$ 与函数 $F_1, F_2$，$\IG_i^{(a F_1 + b F_2)} = a\,\IG_i^{(F_1)} + b\,\IG_i^{(F_2)}$；
  \item \textbf{对称保持性}：若 $F$ 对特征 $i, j$ 对称且 $x_i = x_j,\; x'_i = x'_j$，则 $\IG_i = \IG_j$。
\end{enumerate}
\end{theorem}
\end{thmBox}

\begin{proof}
构造沿着直线段 $\gamma(\alpha) = x' + \alpha(x - x')$ 的一维复合函数 $g(\alpha) = F(\gamma(\alpha))$，其中 $\alpha \in [0, 1]$。容易验证边界取值：
\[
  g(0) = F(x'), \qquad g(1) = F(x), \qquad \gamma'(\alpha) = x - x'.
\]
对 $g(\alpha)$ 几乎处处应用多元链式法则，得
\[
  g'(\alpha) = \frac{d}{d\alpha} F\big(x' + \alpha(x - x')\big)
  = \sum_{i=1}^d \frac{\partial F\big(x' + \alpha(x - x')\big)}{\partial x_i}\,(x_i - x'_i).
\]
在 $[0, 1]$ 上对等式两端积分。由于有限维求和与积分可交换，再结合微积分基本定理（对于 ReLU 等分段光滑网络，有限个折点测度为零，绝对连续函数依然严格适用 Lebesgue 意义下的微积分基本定理）：
\begin{align*}
  F(x) - F(x')
  &= g(1) - g(0)
   = \int_0^1 g'(\alpha)\,d\alpha \\
  &= \int_0^1 \sum_{i=1}^d (x_i - x'_i)\,\frac{\partial F\big(x' + \alpha(x - x')\big)}{\partial x_i}\,d\alpha \\
  &= \sum_{i=1}^d (x_i - x'_i)\int_0^1 \frac{\partial F\big(x' + \alpha(x - x')\big)}{\partial x_i}\,d\alpha
   = \sum_{i=1}^d \IG_i(x; x').
\end{align*}
这就严格证明了完备性恒等式 \eqref{eq:ig-completeness}。

至于四条推论：若 $x$ 与 $x'$ 仅在第 $i$ 维不同，则对所有 $j \neq i$ 有 $x_j - x'_j = 0 \implies \IG_j(x; x') = 0$，代入式 \eqref{eq:ig-completeness} 立刻得 $\IG_i(x; x') = F(x) - F(x') \neq 0$，证得敏感性（a）。实现不变性直接源于偏导数只取决于函数本身的局部行为；线性性由微分与积分的线性性直接得出；对称保持性由对称函数沿对称路径 $\gamma(\alpha)$ 处处满足 $\partial_i F(\gamma(\alpha)) = \partial_j F(\gamma(\alpha))$ 立刻得证。
\end{proof}

回过头再验算例 \ref{ex:grad-fail} 中的饱和反例 $F(x) = 1 - \operatorname{ReLU}(1 - x)$（$x'=0, x=2$）：路径 $\gamma(\alpha) = 2\alpha$ 在 $\alpha \in [0, 1/2)$ 上斜率为 $1$，在 $\alpha \in (1/2, 1]$ 上斜率为 $0$，于是
\[
  \IG(2; 0) = (2 - 0)\left(\int_0^{1/2} 1\,d\alpha + \int_{1/2}^1 0\,d\alpha\right) = 2 \times \frac{1}{2} = 1 = F(2) - F(0).
\]
积分梯度完美穿过了饱和区，将完整的 $1.0$ 输出增量准确归功于特征 $x$！

\subsection{中点黎曼和的 $O(1/m^2)$ 收敛阶与基线选择陷阱}

实际计算中，式 \eqref{eq:ig-def} 的连续积分需通过 $m$ 步数值求积离散化。最常使用的\textbf{中点黎曼和（Midpoint Rule）}在等距子区间 $\big[\frac{k-1}{m}, \frac{k}{m}\big]$ 的中点 $\alpha_k = \frac{k - 0.5}{m}$（$k=1,\dots,m$）处采样：
\begin{equation}\label{eq:ig-midpoint}
  \widehat{\IG}_i^{(m)}(x; x') \;=\; \frac{x_i - x'_i}{m}\sum_{k=1}^m \frac{\partial F\!\left(x' + \frac{k - 0.5}{m}(x - x')\right)}{\partial x_i}.
\end{equation}
为什么工程上强烈推荐使用\textbf{中点法}而非左端点或右端点黎曼和？设被积函数 $h(\alpha) = g'(\alpha)$ 在 $[0,1]$ 上二阶连续可微，子区间宽度 $h_m = 1/m$。在每个中点 $\alpha_k$ 处做二阶泰勒展开：
\[
  h(\alpha) = h(\alpha_k) + h'(\alpha_k)(\alpha - \alpha_k) + \frac{1}{2}h''(\alpha_k)(\alpha - \alpha_k)^2 + O(h_m^3).
\]
在对称区间 $[\alpha_k - h_m/2,\, \alpha_k + h_m/2]$ 上积分时，一阶奇函数项 $\int (\alpha - \alpha_k)\,d\alpha$ \textbf{精确抵消为零}！仅剩二阶项贡献 $\frac{h''(\alpha_k)}{24}h_m^3$。将 $m$ 个子区间的误差累加，全局求积误差严格为
\begin{equation}\label{eq:ig-quad-err}
  \abs{\sum_{i=1}^d \widehat{\IG}_i^{(m)}(x; x') - \big(F(x) - F(x')\big)}
  \;=\; \abs{\frac{1}{24\,m^2}\int_0^1 g'''(\alpha)\,d\alpha + O(m^{-4})}
  \;=\; O\!\left(\frac{1}{m^2}\right).
\end{equation}
相比之下，左/右端点黎曼和的一阶项无法抵消，收敛阶只有缓慢的 $O(1/m)$；而中点法把步数扩大 $10$ 倍，误差便暴降 $100$ 倍！更妙的是，完备性公理本身为我们提供了一个免费的\textbf{在线精度自检探针}：只要运行时实时检查 $\big|\sum_i \widehat{\IG}_i^{(m)} - (F(x) - F(x'))\big|$，一旦误差过大就增大步数 $m$（通常 $m \in [50, 300]$ 即足够精确）。

\begin{insightBox}{基线 $x'$ 的隐藏盲区：为什么不能盲目选全黑图像？}
观察式 \eqref{eq:ig-def} 开头的前置因子 $(x_i - x'_i)$：无论偏导数有多大，\textbf{只要某维特征取值恰好等于基线（$x_i = x'_i$），其积分梯度必定强制为零！}在计算机视觉中，若以全黑图像（像素全 $0$）作为基线去解释一张包含黑色狗或黑色字体的照片，所有黑色像素都会被赋予零归因——模型仿佛“没看见”这条黑狗。工业界与顶会常用\textbf{期望梯度（Expected Gradients）}或高斯模糊基线：从参考分布 $\mathcal{D}'$ 中随机抽取多张背景基线求平均，从而消除单一基线颜色造成的盲区。
\end{insightBox}

\subsection{实测验证：积分梯度完备性误差的 $O(1/m^2)$ 收敛曲线}

配套实验脚本 \texttt{code/vol7\_experiments.py} 中的 \texttt{exp4\_integrated\_\allowbreak gradients\_and\_fairness} 在一个光滑非线性 MLP（$5 \to 16 \to 1$，Tanh 激活，固定随机种子 \texttt{2026}）上，计算输入向量 $x = (1.2,\, -0.8,\, 0.5,\, 1.5,\, -0.3)$ 相对于零基线 $x_0 = \mathbf{0}$ 的真实输出增量 $\Delta F = F(x) - F(x_0)$（其值为 $0.039029$），并扫描求积步数 $m \in \{2, 5, 10, 20, 50, 100, 200\}$。实测数据保存在 \texttt{data/ig\_completeness.csv}，收敛曲线如图 \ref{fig:ig-conv} 所示，明细列于表 \ref{tab:ig-conv}。

\begin{figure}[htbp]
  \centering
  \begin{tikzpicture}
    \begin{loglogaxis}[
      width=0.84\linewidth,
      height=6.0cm,
      xlabel={中点黎曼求积步数 $m$（对数坐标）},
      ylabel={完备性绝对误差 $\big|\sum_i \widehat{\IG}_i^{(m)} - \Delta F\big|$},
      xmin=1.6, xmax=260,
      ymin=3e-9, ymax=1.2e-3,
      grid=both,
      major grid style={dashed, tealLine!75},
      minor grid style={dotted, tealLine!35},
      legend style={at={(0.03,0.06)}, anchor=south west, font=\footnotesize, fill=white, draw=tealLine},
      label style={font=\small\sffamily, color=ink},
      tick label style={font=\footnotesize, color=slate},
      axis line style={color=slate}
    ]
      \addplot[color=teal, very thick, mark=*, mark size=2.3pt]
        table[x=steps, y=abs_err, col sep=comma]{data/ig_completeness.csv};
      \addlegendentry{实测中点法完备性误差（\texttt{data/ig\_completeness.csv}）}

      \addplot[color=ochre, thick, dashed, domain=2:200, samples=30]
        {1.5e-3 / (x^2)};
      \addlegendentry{理论斜率参考线 $C / m^2$（二阶收敛 $O(m^{-2})$）}
    \end{loglogaxis}
  \end{tikzpicture}
  \caption{积分梯度中点黎曼和的完备性误差双对数收敛曲线（$\Delta F = F(x)-F(x_0) = 0.039029$）。在 $m=2 \to 100$ 区间，实测曲线与斜率 $-2$ 的理论二阶收敛参考线高度平行；步数增至 $m=200$ 时误差压低至 $7.45 \times 10^{-9}$，已触及单精度浮点（\texttt{float32}）的舍入噪声底线。}
  \label{fig:ig-conv}
\end{figure}

\begin{table}[htbp]
  \centering
  \caption{积分梯度完备性公理在不同求积步数 $m$ 下的实测误差（$\Delta F = 0.039029$）}
  \label{tab:ig-conv}
  \small
  \begin{tabular}{@{}cccc@{}}
    \toprule
    步数 $m$ & 归因总和 $\sum_{i=1}^5 \widehat{\IG}_i^{(m)}$ & 绝对误差 $|\sum \widehat{\IG}_i^{(m)} - \Delta F|$ & 相对误差 \\
    \midrule
    $2$   & $0.038633$ & $3.95 \times 10^{-4}$ & $1.01\%$ \\
    $5$   & $0.038971$ & $5.78 \times 10^{-5}$ & $0.148\%$ \\
    $10$  & $0.039015$ & $1.42 \times 10^{-5}$ & $0.0365\%$ \\
    $20$  & $0.039025$ & $3.52 \times 10^{-6}$ & $0.0090\%$ \\
    $50$  & $0.039028$ & $5.25 \times 10^{-7}$ & $0.0013\%$ \\
    $100$ & $0.039029$ & $1.19 \times 10^{-7}$ & $0.0003\%$ \\
    $200$ & $0.039029$ & $7.45 \times 10^{-9}$ & $< 0.00002\%$ \\
    \bottomrule
  \end{tabular}
\end{table}

请注意一个极其真实的数值细节：从 $m=5$ 到 $m=50$（步数扩大 $10$ 倍），误差从 $5.78 \times 10^{-5}$ 精确降至 $5.25 \times 10^{-7}$（缩小约 $110$ 倍，完全吻合 $O(1/m^2)$）；而到了 $m=200$，误差跌破 $10^{-8}$（恰为 $2^{-27} \approx 7.45 \times 10^{-9}$），已进入 PyTorch 默认 \texttt{float32} 尾数 24-bit 精度的舍入底噪区。下面是带有自动完备性断言的干净 PyTorch 实现：

\begin{lstlisting}[caption={带完备性公理自检的积分梯度（中点黎曼求积，PyTorch）}]
import torch

def integrated_gradients(model, x, x_base, steps: int = 50, tol: float = 1e-4):
    """Compute midpoint-rule IG for a single sample x (shape [1, d])."""
    # Midpoints alpha_k = (k - 0.5) / steps in (0, 1) -> O(1/m^2) accuracy
    alphas = torch.linspace(0.5 / steps, 1.0 - 0.5 / steps, steps, device=x.device)
    alphas = alphas.view(-1, *([1] * (x.ndim - 1)))
    path_pts = (x_base + alphas * (x - x_base)).detach().requires_grad_(True)
    # Single batched forward+backward across all interpolation steps
    out = model(path_pts).sum()
    grads = torch.autograd.grad(out, path_pts)[0]
    ig = (x - x_base) * grads.mean(dim=0, keepdim=True)
    # Self-verify Completeness Axiom: sum(IG) == F(x) - F(x_base)
    with torch.no_grad():
        delta_f = float((model(x) - model(x_base)).item())
    err = abs(float(ig.sum().item()) - delta_f)
    assert err < tol, f"Completeness error {err:.2e} exceeds tol={tol}; increase steps!"
    return ig, err
\end{lstlisting}

\section{合作博弈论视角：Shapley 值与 SHAP 特征归因}

积分梯度依赖模型对输入可微；但如果我们要解释一个不可微的随机森林、GBDT（第 4 卷第 5 章）或离散黑箱 API，该如何公平地把预测收益分摊给每个特征？1953 年，经济学家 Lloyd Shapley（2012 年诺贝尔经济学奖得主）在合作博弈论中给出了唯一答案；Lundberg 与 Lee（NeurIPS 2017）将其统一引入机器学习，提出了 \textbf{SHAP（SHapley Additive exPlanations）}框架。

\subsection{Shapley 唯一性定理与四大公理}

设特征全集为 $N = \{1, 2, \dots, d\}$，每个特征子集（联盟）$S \subseteq N$ 对应一个特征价值函数 $v: 2^N \to \R$（在机器学习中通常取条件期望 $v(S) = \E[F(X) \mid X_S = x_S]$，并记空联盟基线 $v(\emptyset) = \E[F(X)]$）。当特征 $i \notin S$ 加入联盟 $S$ 时，它带来的\textbf{边际贡献（Marginal Contribution）}为 $v(S \cup \{i\}) - v(S)$。

\begin{thmBox}
\begin{theorem}[Shapley 唯一性定理]\label{thm:shapley}
在所有把联盟价值 $v(N) - v(\emptyset)$ 分配给 $d$ 个参与者的方案 $\phi(v) = (\phi_1(v), \dots, \phi_d(v))^\top$ 中，\textbf{存在且唯一存在}一种分配方案同时满足以下四条公理：
\begin{enumerate}
  \item \textbf{效率性（Efficiency）}：$\sum_{i=1}^d \phi_i(v) = v(N) - v(\emptyset)$；
  \item \textbf{对称性（Symmetry）}：若对所有 $S \subseteq N \setminus \{i, j\}$ 均有 $v(S \cup \{i\}) = v(S \cup \{j\})$，则 $\phi_i(v) = \phi_j(v)$；
  \item \textbf{哑元性（Dummy / Null Player）}：若对所有 $S \subseteq N \setminus \{i\}$ 均有 $v(S \cup \{i\}) = v(S)$，则 $\phi_i(v) = 0$；
  \item \textbf{可加性（Additivity）}：对任意两个博弈 $v, w$，有 $\phi_i(v + w) = \phi_i(v) + \phi_i(w)$。
\end{enumerate}
该唯一分配即为 \textbf{Shapley 值}：
\begin{equation}\label{eq:shapley-formula}
  \phi_i(v) \;=\; \sum_{S \subseteq N \setminus \{i\}} \frac{|S|!\,(d - |S| - 1)!}{d!}\,\Big[v(S \cup \{i\}) - v(S)\Big]
  \;=\; \frac{1}{d!}\sum_{\pi \in \mathfrak{S}_d} \Big[v\big(P_i^\pi \cup \{i\}\big) - v\big(P_i^\pi\big)\Big],
\end{equation}
其中 $\mathfrak{S}_d$ 是 $N$ 上全部 $d!$ 个全排列的集合，$P_i^\pi$ 表示在排列 $\pi$ 中排在特征 $i$ 之前的所有特征集合。
\end{theorem}
\end{thmBox}

式 \eqref{eq:shapley-formula} 的组合权重极其自然：想象 $d$ 个特征以完全随机的顺序 $\pi$ 依次“进场”，特征 $i$ 进场时恰好看到已有联盟 $S$ 在场的排列数等于 $|S|!$（$S$ 内部任意排）乘以 $(d - |S| - 1)!$（剩余成员在 $i$ 之后任意排），除以总排列数 $d!$ 即得式中的权重。两个表达式由此严格等价。此外，注意到对任意排列 $\pi$，依次累加每位进场者的边际贡献构成一个精确的裂项求和（Telescoping Sum）：
\[
  \sum_{i=1}^d \Big[v\big(P_i^\pi \cup \{i\}\big) - v\big(P_i^\pi\big)\Big] = v(N) - v(\emptyset),
\]
再对所有 $d!$ 个排列取平均，效率性公理立即迎刃而解！对称性、哑元性与可加性亦可直接从公式的线性对称结构一眼看出。唯一性的标准证明则通过构造一致博弈（Unanimity Games）$u_T(S) = \mathbb{I}[T \subseteq S]$（$\emptyset \neq T \subseteq N$）这组 $2^d - 1$ 维向量空间基：由效率、对称、哑元三条公理唯一钉死 $\phi_i(c\,u_T) = c/|T|$（若 $i \in T$）或 $0$（若 $i \notin T$），再由可加性线性延拓到全部博弈空间即证唯一。

\begin{example}[三特征合作博弈手算：协同交互与四公理全面核验]\label{ex:shapley-3feat}
考虑一个 $d = 3$ 特征的贷款风控评分博弈 $N = \{1, 2, 3\}$，其中特征 $1$（稳定收入）与特征 $2$（良好征信）存在显著正向协同，而特征 $3$（手机壁纸颜色）为纯噪声特征。设各联盟价值（单位：分）如下：
\[
  v(\emptyset) = 0, \quad
  v(\{1\}) = 2, \quad v(\{2\}) = 2, \quad v(\{3\}) = 0,
\]
\[
  v(\{1,2\}) = 8, \quad v(\{1,3\}) = 2, \quad v(\{2,3\}) = 2, \quad v(\{1,2,3\}) = 8.
\]
我们列出全部 $3! = 6$ 种特征入场排列，逐一记录每个特征入场时的边际贡献：

\begin{center}
\small
\begin{tabular}{@{}lcccc@{}}
  \toprule
  排列 $\pi$ & 特征 $1$ 边际贡献 & 特征 $2$ 边际贡献 & 特征 $3$ 边际贡献 & 排列总和 \\
  \midrule
  $(1, 2, 3)$ & $v(\{1\})-v(\emptyset)=2$ & $v(\{1,2\})-v(\{1\})=6$ & $v(\{1,2,3\})-v(\{1,2\})=0$ & $8$ \\
  $(1, 3, 2)$ & $v(\{1\})-v(\emptyset)=2$ & $v(\{1,2,3\})-v(\{1,3\})=6$ & $v(\{1,3\})-v(\{1\})=0$ & $8$ \\
  $(2, 1, 3)$ & $v(\{1,2\})-v(\{2\})=6$ & $v(\{2\})-v(\emptyset)=2$ & $v(\{1,2,3\})-v(\{1,2\})=0$ & $8$ \\
  $(2, 3, 1)$ & $v(\{1,2,3\})-v(\{2,3\})=6$ & $v(\{2\})-v(\emptyset)=2$ & $v(\{2,3\})-v(\{2\})=0$ & $8$ \\
  $(3, 1, 2)$ & $v(\{1,3\})-v(\{3\})=2$ & $v(\{1,2,3\})-v(\{1,3\})=6$ & $v(\{3\})-v(\emptyset)=0$ & $8$ \\
  $(3, 2, 1)$ & $v(\{1,2,3\})-v(\{2,3\})=6$ & $v(\{2,3\})-v(\{3\})=2$ & $v(\{3\})-v(\emptyset)=0$ & $8$ \\
  \midrule
  \textbf{均值 $\phi_i(v)$} & $(2+2+6+6+2+6)/6 = \mathbf{4}$ & $(6+6+2+2+6+2)/6 = \mathbf{4}$ & $\mathbf{0}$ & $\mathbf{8}$ \\
  \bottomrule
\end{tabular}
\end{center}

逐条核验四公理：\textbf{效率性} $\phi_1 + \phi_2 + \phi_3 = 4 + 4 + 0 = 8 = v(N) - v(\emptyset)$（不仅分掉了两人各自的 $2$ 分，还把协同产生的 $4$ 分冗余收益公平平分！）；\textbf{对称性} 特征 $1, 2$ 角色完全对称，果真 $\phi_1 = \phi_2 = 4$；\textbf{哑元性} 噪声特征 $3$ 严格得 $\phi_3 = 0$；\textbf{可加性} 若将本博弈拆成线性基础分 $v_{\text{lin}}(S) = 2|S \cap \{1,2\}|$ 与纯协同项 $v_{\text{syn}}(S) = 4\,\mathbb{I}[\{1,2\}\subseteq S]$，则 $\phi(v_{\text{lin}})=(2,2,0)$、$\phi(v_{\text{syn}})=(2,2,0)$，相加恰好等于 $(4,4,0)$！
\end{example}

\subsection{指数复杂度壁垒与 KernelSHAP / TreeSHAP 近似}

虽然 Shapley 值在公理上无懈可击，但直接套用式 \eqref{eq:shapley-formula} 需要遍历 $2^d$ 个子集。当特征维数 $d = 100$ 时，$2^{100} \approx 1.27 \times 10^{30}$ 次模型前向评估在宇宙寿命内也算不完。工程界为此演化出两大主力算法：
\begin{itemize}
  \item \textbf{KernelSHAP（模型无关加权回归）}：Lundberg 与 Lee 证明，Shapley 值等价于一个特殊的加权线性最小二乘回归解！具体而言，随机采样 $M \ll 2^d$ 个二值特征掩码 $z' \in \{0, 1\}^d$（去掉全 $0$ 与全 $1$），构造解释线性模型 $g(z') = \phi_0 + \sum_{i=1}^d \phi_i z'_i$，赋予 \textbf{Shapley 核权重}
  \[
    \pi_x(z') = \frac{d - 1}{\binom{d}{|z'|}\,|z'|\,(d - |z'|)},
  \]
  并在约束 $\phi_0 = v(\emptyset)$、$\sum_{i=1}^d \phi_i = v(N) - v(\emptyset)$ 下最小化加权平方误差 $\sum_{z'} \pi_x(z')\big(v(z') - g(z')\big)^2$。核权重对极小联盟（$|z'|=1$）与极大联盟（$|z'|=d-1$）赋予最高权重，恰恰因为这两端能最清晰地分离单个特征的边际效应。
  \item \textbf{TreeSHAP（树模型多项式精确算法）}：对于 XGBoost、LightGBM 等决策树集成，树结构天然将特征交互限制在单条根到叶的路径上。Lundberg 等人（Nature Machine Intelligence 2020）通过沿树深度动态规划追踪落在子集 $S$ 内的特征比例，将精确 Shapley 值的计算复杂度从 $O(T L \cdot 2^d)$ 骤降至多项式时间 $O(T L D^2)$（$T$ 为树棵数，$L$ 为最大叶节点数，$D$ 为最大树深），实现了百万级表格样本的毫秒级精确归因。
  \item \textbf{与积分梯度的深刻统一}：如果把连续特征向量 $x \in \R^d$ 视为连续合作博弈（Aumann--Shapley 理论），积分梯度恰恰是连续博弈中的 \textbf{Aumann--Shapley 路径值}——两者在离散与连续两个世界遥相呼应。
\end{itemize}

\section{机械可解释性：探针、激活修补与稀疏自编码器（SAE）}

输入--输出归因（如 IG 与 SHAP）告诉我们“哪几个输入词或像素影响了最终输出”，却仍然把神经网络内部成百上千层当作不可见的黑盒。近年来兴起的\textbf{机械可解释性（Mechanistic Interpretability）}则提出了一个更具雄心的目标：像逆向工程二进制汇编代码那样，将神经网络的权重与隐藏激活彻底拆解为人类可读的\textbf{概念特征（Features）}与\textbf{因果电路（Circuits）}。

\subsection{线性探针的诱惑与“选择性”陷阱}

最古老的白箱检测工具是\textbf{线性探针（Linear Probing）}：冻结预训练模型，提取第 $\ell$ 层隐藏激活 $h^{(\ell)}(x) \in \R^{d_{\text{model}}}$，再训练一个简单线性分类器 $\hat{y}_{\text{prop}} = \softmax(W_p h^{(\ell)}(x) + b_p)$ 来预测某个语言学或语义属性（如词性、时态、棋盘局面或事实真假）。如果探针准确率高达 $92\%$，是否就证明模型在第 $\ell$ 层“学懂了并用这个概念做推理”？

\textbf{绝非如此！}Hewitt 与 Liang（EMNLP 2019）揭示了著名的\textbf{探针记忆幻觉}：当隐藏维数 $d_{\text{model}} = 4096$ 足够高时，即使把输入特征做随机正交投影，高维空间中的词表嵌入依然近似线性可分——探针完全可以靠“死记硬背每个词的 ID”就在标注属性上拿到高分！因此，必须报告\textbf{选择性（Selectivity）}：
\[
  \mathrm{Selectivity} \;=\; \mathrm{Acc}_{\text{linguistic}} - \mathrm{Acc}_{\text{control}},
\]
其中 $\mathrm{Acc}_{\text{control}}$ 是探针在一个\textbf{随机打乱词级标签的对照任务（Control Task）}上的准确率。只有当真实任务准确率远高于对照任务（高选择性），才能初步相信该表示线性编码了结构化概念；即便如此，探针仍只能证明\textbf{相关性（信息在那里）}，无法证明\textbf{因果性（下游层真的读取并使用了它）}。

\subsection{激活修补（Activation Patching）：把因果干预请进 Transformer}

要从“相关”跨越到“因果”，必须引入 Judea Pearl 因果推断中的 $\mathrm{do}$-干预算子——在深度学习中，这就是 \textbf{激活修补（Activation Patching / Causal Tracing，Meng et al., NeurIPS 2022）}。以经典事实召回提示为例：干净输入 $x_{\text{clean}} = \text{``埃菲尔铁塔位于''}$，模型高概率输出 $\text{``巴黎''}$。我们执行三趟前向传播：
\begin{enumerate}
  \item \textbf{干净前向（Clean Run）}：输入 $x_{\text{clean}}$，缓存每一层、每个位置上的残差流或子模块激活 $h^{(\ell, t)}_{\text{clean}}$，记录高 logit 差 $L_{\text{clean}} = \mathrm{logit}(\text{巴黎}) - \mathrm{logit}(\text{罗马})$；
  \item \textbf{损坏前向（Corrupted Run）}：对主语词元（``埃菲尔铁塔''）的输入嵌入加入高斯噪声（或替换为另一主语 ``斗兽场''），得到损坏输入 $x_{\text{corr}}$，此时输出崩溃，logit 差跌至低值 $L_{\text{corr}}$；
  \item \textbf{单点修补前向（Patched Run）}：在运行损坏前向的过程中，\textbf{仅将特定层 $\ell$、特定词元位置 $t$ 的隐藏状态强制覆写为干净缓存值} $\mathrm{do}\!\big(h^{(\ell, t)} \leftarrow h^{(\ell, t)}_{\text{clean}}\big)$，其余一切保持损坏状态，观测输出 logit 差恢复到的水平 $L_{\text{patch}}(\ell, t)$。
\end{enumerate}
据此定义该节点对事实召回的\textbf{标准化间接因果效应（Normalized Indirect Effect）}：
\begin{equation}\label{eq:activation-patching}
  \mathrm{IE}(\ell, t) \;=\; \frac{L_{\text{patch}}(\ell, t) - L_{\text{corr}}}{L_{\text{clean}} - L_{\text{corr}}}.
\end{equation}
如果仅仅把中层某个 MLP（例如主语最后一个词元位置的第 $\ell^*$ 层 MLP）换回干净激活，模型输出就立刻从“罗马”大举跳回“巴黎”（$\mathrm{IE} \approx 0.7$），我们就获得了铁一般的\textbf{因果证据}：该模块是这一推理电路中不可或缺的因果枢纽！图 \ref{fig:ig-and-patching} 将积分梯度的连续路径几何与激活修补的离散因果干预并置展示。

\begin{figure}[htbp]
  \centering
  \resizebox{0.98\linewidth}{!}{%
  \begin{tikzpicture}[
    >=Stealth,
    box/.style={rectangle, rounded corners=1.4mm, draw=teal, thick, fill=tealLight,
      minimum width=16mm, minimum height=6.5mm, align=center, font=\sffamily\scriptsize, text=ink},
    cbox/.style={rectangle, rounded corners=1.4mm, draw=ochre, thick, fill=ochreLight,
      minimum width=16mm, minimum height=6.5mm, align=center, font=\sffamily\scriptsize, text=ink},
    pbox/.style={rectangle, rounded corners=1.4mm, draw=navy, very thick, fill=navyLight,
      minimum width=18mm, minimum height=6.5mm, align=center, font=\sffamily\scriptsize, text=navy}
  ]
    % Left Panel: Integrated Gradients Straight-Line Path
    \begin{scope}[shift={(0,0)}]
      \draw[rounded corners=2mm, draw=tealLine, fill=white] (-0.4,-0.3) rectangle (6.5,5.3);
      \node[font=\sffamily\bfseries\footnotesize, text=tealDark] at (3.05, 4.92) {(a) 积分梯度：穿越饱和区的直线路径 $\gamma(\alpha)$};

      % Saturation region shading
      \fill[ochreLight!85] (2.9, 0.5) rectangle (6.1, 4.35);
      \draw[ochre, dashed, thick] (2.9, 0.5) -- (2.9, 4.35);
      \node[font=\sffamily\scriptsize, text=ochre, align=center] at (4.65, 3.88) {饱和平坦区\\$\nabla_x F \approx \mathbf{0}$};
      \node[font=\sffamily\scriptsize, text=tealDark, align=center] at (1.35, 3.88) {活跃斜坡区\\$\nabla_x F \neq \mathbf{0}$};

      % Axes
      \draw[->, slate, thick] (0.2, 0.5) -- (6.15, 0.5) node[right, font=\scriptsize] {$x_1$};
      \draw[->, slate, thick] (0.2, 0.5) -- (0.2, 4.45) node[above, font=\scriptsize] {$x_2$};

      % Baseline x' and Input x
      \coordinate (x0) at (0.85, 1.15);
      \coordinate (x1) at (5.35, 3.10);
      \draw[tealDark, very thick] (x0) -- (x1);

      % Midpoint quadrature dots & gradient ticks
      \foreach \a in {0.1, 0.3, 0.5, 0.7, 0.9} {
        \coordinate (p) at ($(x0)!\a!(x1)$);
        \fill[gold] (p) circle (2.2pt);
      }
      % Gradient arrows in active zone vs zero at x
      \draw[->, teal, thick] ($(x0)!0.1!(x1)$) -- ++(0.32, 0.42);
      \draw[->, teal, thick] ($(x0)!0.3!(x1)$) -- ++(0.32, 0.42);
      \node[font=\scriptsize\sffamily, text=teal, right] at (1.28, 2.28) {沿途累加梯度};

      \fill[navy] (x0) circle (3pt) node[below right, font=\scriptsize\sffamily, text=navy] {基线 $x'$};
      \fill[ochre] (x1) circle (3pt);
      \node[font=\scriptsize\sffamily, text=ochre, align=center] at (4.95, 2.38) {输入 $x$\\（局部梯度为 $0$）};
      \node[font=\scriptsize\sffamily, text=tealDark, rotate=23] at (2.85, 1.62) {$\gamma(\alpha) = x' + \alpha(x - x')$};
    \end{scope}

    % Right Panel: Activation Patching Causal Tracing
    \begin{scope}[shift={(7.1,0)}]
      \draw[rounded corners=2mm, draw=tealLine, fill=white] (-0.3,-0.3) rectangle (7.4,5.3);
      \node[font=\sffamily\bfseries\footnotesize, text=tealDark] at (3.55, 4.92) {(b) 激活修补：隔离 Transformer 因果电路};

      % Clean run column
      \node[font=\sffamily\bfseries\scriptsize, text=tealDark] at (1.15, 4.44) {干净前向};
      \node[box] (c0) at (1.15, 3.82) {$x_{\text{clean}}$ 嵌入};
      \node[box] (c1) at (1.15, 2.45) {层 $\ell$ 残差 $h^{(\ell)}_{\text{clean}}$};
      \node[box] (c2) at (1.15, 1.32) {高层 $\ell+1\dots L$};
      \node[font=\sffamily\scriptsize, text=tealDark] (cout) at (1.15, 0.35) {输出: ``巴黎'' ($+$)};
      \draw[->, teal, thick] (c0) -- (c1);
      \draw[->, teal, thick] (c1) -- (c2);
      \draw[->, teal, thick] (c2) -- (cout);

      % Corrupted run column
      \node[font=\sffamily\bfseries\scriptsize, text=ochre] at (3.65, 4.44) {损坏前向};
      \node[cbox] (r0) at (3.65, 3.82) {$x_{\text{corr}}$ 噪声主语};
      \node[cbox] (r1) at (3.65, 2.45) {层 $\ell$ 残差 $h^{(\ell)}_{\text{corr}}$};
      \node[cbox] (r2) at (3.65, 1.32) {高层 $\ell+1\dots L$};
      \node[font=\sffamily\scriptsize, text=ochre] (rout) at (3.65, 0.35) {输出崩溃 ($-$)};
      \draw[->, ochre, thick] (r0) -- (r1);
      \draw[->, ochre, thick] (r1) -- (r2);
      \draw[->, ochre, thick] (r2) -- (rout);

      % Patched run column
      \node[font=\sffamily\bfseries\scriptsize, text=navy] at (6.15, 4.44) {修补干预前向};
      \node[cbox] (p0) at (6.15, 3.82) {$x_{\text{corr}}$ 噪声主语};
      \node[pbox] (p1) at (6.15, 2.45) {$\mathrm{do}(h^{(\ell)}\!\leftarrow\!h^{(\ell)}_{\text{clean}})$};
      \node[box] (p2) at (6.15, 1.32) {高层 $\ell+1\dots L$};
      \node[font=\sffamily\bfseries\scriptsize, text=navy] (pout) at (6.15, 0.35) {恢复 ``巴黎'' ($\mathrm{IE}\!\uparrow$)};
      \draw[->, ochre, thick] (p0) -- (p1);
      \draw[->, navy, very thick] (p1) -- (p2);
      \draw[->, navy, very thick] (p2) -- (pout);

      % Patching curved arrow from clean h^(l) to patched h^(l)
      \draw[->, gold, very thick, dashed] (c1.north east) .. controls (2.8, 3.22) and (4.5, 3.22) .. (p1.north west);
      \node[font=\sffamily\bfseries\scriptsize, text=gold!90!black, fill=white, inner sep=1.5pt]
        at (3.65, 3.16) {精准注入干净激活};
    \end{scope}
  \end{tikzpicture}}
  \caption{两种核心可解释性范式的几何与因果对比。(a) 积分梯度沿直线路径累积梯度，跨越饱和平坦区，满足完备性公理；(b) 激活修补通过 $\mathrm{do}$-干预将干净残差流精准植入损坏前向，用输出恢复度 $\mathrm{IE}(\ell, t)$ 锁定真正起作用的因果电路。}
  \label{fig:ig-and-patching}
\end{figure}

\subsection{神经叠加（Superposition）与稀疏自编码器（SAE）}

即便通过激活修补锁定了某关键 MLP 层，直接观察单个神经元仍常令人困惑：同一个神经元竟然同时对“猫的耳朵”“汽车前灯”和“Python 缩进”产生高激活！这种现象称为\textbf{多义性（Polysemanticity）}。Elhage 等人（Anthropic, 2022）提出了极具解释力的\textbf{叠加假说（Superposition Hypothesis）}：现实世界需要表征的独立概念数目 $M$ 远远大于神经网络的隐藏维数 $d$（$M \gg d$）；但幸运的是，自然输入中的绝大多数概念是\textbf{高度稀疏}的（同一时刻只有极少数概念被激活）。根据高维几何中的 Johnson--Lindenstrauss 引理，在 $\R^d$ 空间中可以塞下指数级数量（$e^{\Omega(\varepsilon^2 d)}$）个彼此几乎正交（夹角余弦 $|\cos\theta| < \varepsilon$）的单位方向！于是神经网络把 $M \gg d$ 个概念“叠加”压缩进 $d$ 维残差流中，以微小的线性串扰换取巨大的表征容量。

如何把叠加在一起的多义激活 $h \in \R^d$ 重新拆解回干净的单义概念？答案直接回连到本系列\textbf{第 4 卷第 2 章的 Lasso 与稀疏字典学习}——训练一个\textbf{过完备稀疏自编码器（Sparse Autoencoder, SAE）}！设隐空间维度 $M \gg d$（例如 $M = 16 d$ 或 $64 d$），定义编码器与解码器：
\begin{equation}\label{eq:sae}
  z = \operatorname{ReLU}\!\big(W_{\text{enc}}(h - b_{\text{dec}}) + b_{\text{enc}}\big) \in \R_{\ge 0}^M,
  \qquad
  \hat{h} = W_{\text{dec}} z + b_{\text{dec}} \in \R^d,
\end{equation}
通过最小化重构误差与 $L_1$ 稀疏惩罚之和来学习过完备字典 $W_{\text{dec}} \in \R^{d \times M}$：
\begin{equation}\label{eq:sae-loss}
  \mathcal{L}_{\text{SAE}} \;=\; \E_{h}\!\Big[\norm{h - \hat{h}}_2^2 \;+\; \lambda \norm{z}_1\Big],
  \qquad \text{s.t. } \norm{W_{\text{dec}, :, j}}_2 = 1 \;\; (\forall j = 1,\dots,M).
\end{equation}
请留意式 \eqref{eq:sae-loss} 右侧那条不可或缺的\textbf{单位列范数约束}：如果不约束解码器字典向量的模长，优化器就会作弊——把 $W_{\text{dec}}$ 无限放大、把 $z$ 无限缩小，在完全不改变重构 $\hat{h}$ 的前提下把 $\lambda\norm{z}_1$ 压到趋近于 $0$！在单位范数约束下，$L_1$ 惩罚迫使每个样本只激活极少数（如 $L_0 \approx 20 \sim 50$ 个）字典原子。2024 年前后对前沿大语言模型残差流的大规模 SAE 分解，成功提取出了“金门大桥”“代码安全漏洞”“奉承语气”等高度单义的特征方向；甚至人为调大某个 SAE 特征的激活幅值（Feature Steering），就能精确定向操控模型的生成风格与安全拒绝行为。

\section{算法公平性数学框架与 Chouldechova 不可能定理}

可解释性让我们看清模型在想什么，而\textbf{算法公平性（Algorithmic Fairness）}追问的是：当模型决策关乎贷款额度、招聘筛选、疾病分诊或保释评估时，什么是数学上“公平”的预测？这一领域最具震撼力的里程碑发现是：\textbf{人类直觉中几条同样天经地义的公平定义，在数学上是彼此水火不容的！}

\subsection{群体公平性的三大核心定义}

考虑二分类决策场景：个体特征为 $X$，受保护群体属性（如性别、地域、年龄段等）为 $A \in \{0, 1\}$，真实二元结果（如是否按期还款或是否再犯）为 $Y \in \{0, 1\}$，模型给出的预测概率为 $S \in [0, 1]$，二元决策为 $\hat{Y} \in \{0, 1\}$。对每个群体 $g \in \{0, 1\}$，定义其混淆矩阵核心指标：
\begin{itemize}
  \item \textbf{群体基础率（Base Rate）}：$\pi_g = \Prob(Y = 1 \mid A = g)$；
  \item \textbf{假阳性率（误报率，FPR）}：$\mathrm{FPR}_g = \Prob(\hat{Y} = 1 \mid Y = 0,\, A = g)$；
  \item \textbf{真阳性率（TPR）与假阴性率（漏报率，FNR）}：$\mathrm{TPR}_g = \Prob(\hat{Y} = 1 \mid Y = 1,\, A = g)$，$\mathrm{FNR}_g = 1 - \mathrm{TPR}_g$；
  \item \textbf{阳性预测值（精确率，PPV）}：$\mathrm{PPV}_g = \Prob(Y = 1 \mid \hat{Y} = 1,\, A = g)$。
\end{itemize}

\begin{defBox}
\begin{definition}[群体公平性的三大经典准则]\label{def:fairness-criteria}
记两组受保护属性为 $g \in \{0, 1\}$：
\begin{enumerate}
  \item \textbf{统计均等（Demographic Parity / Statistical Parity）}：决策结果与群体属性统计独立，即各组获得正决策的比例相等：
  \[
    \Prob(\hat{Y} = 1 \mid A = 0) \;=\; \Prob(\hat{Y} = 1 \mid A = 1).
  \]
  \item \textbf{均等赔率（Equalized Odds，Hardt et al., NeurIPS 2016）}：在给定真实结果 $Y$ 的条件下，决策与群体属性条件独立，即两组拥有\textbf{相同的误报率与漏报率}：
  \[
    \mathrm{FPR}_0 = \mathrm{FPR}_1
    \quad \text{且} \quad
    \mathrm{FNR}_0 = \mathrm{FNR}_1 \;\; (\Longleftrightarrow \mathrm{TPR}_0 = \mathrm{TPR}_1).
  \]
  \item \textbf{预测均等 / 组内校准（Predictive Parity / Calibration within Groups）}：给出正预测后，真实为正的概率跨组一致（分值对所有群体含义相同）：
  \[
    \mathrm{PPV}_0 \;=\; \mathrm{PPV}_1 \qquad \big(\text{连续分值版本：}\;\Prob(Y=1 \mid S=s,\, A=g) = s\big).
  \]
\end{enumerate}
\end{definition}
\end{defBox}

乍看之下，\textbf{组内校准（预测均等）}与\textbf{均等赔率}似乎都理所当然：前者保证“同样被模型判为高风险的两个人，无论来自哪个群体，真实风险概率相同”；后者保证“同样清白的无辜者，不会因为所属群体不同而承受更高的误判率”。然而，2016 年美国司法风控软件 COMPAS 爆发的著名争论（ProPublica 指控其对黑人被告的误报率 $\mathrm{FPR}$ 几乎是白人被告的两倍，而开发商 Northpointe 反驳称其风险评分在两个群体上同样良好校准）促使统计学家 Alexandra Chouldechova 与计算机科学家 Jon Kleinberg 等人独立证明了一条震动整个 AI 伦理界的不可能定理。

\subsection{Chouldechova 公平性不可能定理（完整贝叶斯证明）}

\begin{thmBox}
\begin{theorem}[算法公平性不可能定理，Chouldechova 2017]\label{thm:chouldechova}
对任意群体 $g \in \{0, 1\}$，设群体基础率 $\pi_g \in (0, 1)$，且分类器在该群体上给出非零正预测概率，满足阳性预测值 $\mathrm{PPV}_g \in (0, 1)$。则以下精确的贝叶斯代数恒等式恒成立：
\begin{equation}\label{eq:chouldechova-id}
  \mathrm{FPR}_g \;=\; \frac{\pi_g}{1 - \pi_g}\cdot \frac{1 - \mathrm{PPV}_g}{\mathrm{PPV}_g}\cdot \big(1 - \mathrm{FNR}_g\big).
\end{equation}
\textbf{不可能推论}：只要两个群体的真实基础率不相等（$\pi_0 \neq \pi_1$），那么任何满足\textbf{预测均等}（$\mathrm{PPV}_0 = \mathrm{PPV}_1 = p \in (0, 1)$）的不完美分类器，\textbf{绝对不可能}同时满足\textbf{均等赔率}（$\mathrm{FPR}_0 = \mathrm{FPR}_1$ 且 $\mathrm{FNR}_0 = \mathrm{FNR}_1$）！
\end{theorem}
\end{thmBox}

\begin{proof}
固定群体 $A = g$，为了记号简洁，以下所有概率默认以 $A = g$ 为条件。考虑联合事件 $\{Y = 0,\, \hat{Y} = 1\}$（假阳性）与 $\{Y = 1,\, \hat{Y} = 1\}$（真阳性）。按条件概率乘法公式将两者的联合概率分别按“先 $Y$ 后 $\hat{Y}$”和“先 $\hat{Y}$ 后 $Y$”两种顺序展开：
\begin{align}
  \Prob(Y = 0,\, \hat{Y} = 1)
  &= \Prob(Y = 0)\,\Prob(\hat{Y} = 1 \mid Y = 0) = (1 - \pi_g)\,\mathrm{FPR}_g \label{eq:fp-way1}\\
  &= \Prob(\hat{Y} = 1)\,\Prob(Y = 0 \mid \hat{Y} = 1) = \Prob(\hat{Y} = 1)\,(1 - \mathrm{PPV}_g), \label{eq:fp-way2}
\end{align}
以及
\begin{align}
  \Prob(Y = 1,\, \hat{Y} = 1)
  &= \Prob(Y = 1)\,\Prob(\hat{Y} = 1 \mid Y = 1) = \pi_g\,(1 - \mathrm{FNR}_g) \label{eq:tp-way1}\\
  &= \Prob(\hat{Y} = 1)\,\Prob(Y = 1 \mid \hat{Y} = 1) = \Prob(\hat{Y} = 1)\,\mathrm{PPV}_g. \label{eq:tp-way2}
\end{align}
因为 $\mathrm{PPV}_g > 0$ 且 $\Prob(\hat{Y} = 1) > 0$，由式 \eqref{eq:tp-way1}--\eqref{eq:tp-way2} 可知 $\Prob(Y = 1,\, \hat{Y} = 1) > 0$，且立刻解出 $1 - \mathrm{FNR}_g > 0$。将式 \eqref{eq:fp-way1}--\eqref{eq:fp-way2} 除以式 \eqref{eq:tp-way1}--\eqref{eq:tp-way2}，消去未知的整体阳性决策率 $\Prob(\hat{Y} = 1)$：
\[
  \frac{(1 - \pi_g)\,\mathrm{FPR}_g}{\pi_g\,(1 - \mathrm{FNR}_g)}
  \;=\; \frac{\Prob(\hat{Y} = 1)\,(1 - \mathrm{PPV}_g)}{\Prob(\hat{Y} = 1)\,\mathrm{PPV}_g}
  \;=\; \frac{1 - \mathrm{PPV}_g}{\mathrm{PPV}_g}.
\]
将左端系数 $\frac{1 - \pi_g}{\pi_g\,(1 - \mathrm{FNR}_g)}$ 移到右端，立即得到精确恒等式：
\[
  \mathrm{FPR}_g \;=\; \frac{\pi_g}{1 - \pi_g}\cdot \frac{1 - \mathrm{PPV}_g}{\mathrm{PPV}_g}\cdot \big(1 - \mathrm{FNR}_g\big).
\]
现在用反证法证明不可能推论。假设在 $\pi_0 \neq \pi_1$ 的条件下，存在分类器同时满足 $\mathrm{PPV}_0 = \mathrm{PPV}_1 = p \in (0, 1)$、$\mathrm{FPR}_0 = \mathrm{FPR}_1$ 与 $\mathrm{FNR}_0 = \mathrm{FNR}_1$。因为 $p \in (0, 1)$，有 $\frac{1 - p}{p} > 0$ 且前述已证 $1 - \mathrm{FNR}_g > 0$。将恒等式 \eqref{eq:chouldechova-id} 对 $g=0$ 与 $g=1$ 两边相除（或直接消去相同的非零因子 $\frac{1-p}{p}(1 - \mathrm{FNR}_g)$），必然推出
\[
  \frac{\pi_0}{1 - \pi_0} \;=\; \frac{\pi_1}{1 - \pi_1}
  \quad\Longleftrightarrow\quad
  \pi_0 = \pi_1,
\]
这与前提 $\pi_0 \neq \pi_1$ 直接矛盾！仅有两类极端平凡情形可绕过矛盾：（i）\textbf{完美神谕分类器}（$p = 1$ 且 $\mathrm{FPR}_g = \mathrm{FNR}_g = 0$）；或（ii）\textbf{两组基础率原本完全相等}（$\pi_0 = \pi_1$）。除此之外，三者绝无可能同时满足。证毕。
\end{proof}

\begin{insightBox}{天书般的简洁，社会学的震撼}
请注意式 \eqref{eq:chouldechova-id} 的证明只用了四行初等条件概率——它不依赖任何模型架构，无论是逻辑回归、随机森林还是万亿参数大模型，只要预测不是 $100\%$ 完美且两组基础率不同，这条代数锁链便不可违抗！更重要的是，如果模型维持跨组相同的 $\mathrm{PPV}$，那么\textbf{基础率更高的群体在统计上必然承受更高的误报率 $\mathrm{FPR}$、更低的漏报率 $\mathrm{FNR}$}——这精确重现了 COMPAS 争议的全部定量结构。这也呼应了本章扉页康德的箴言：公平权衡不是一道可以外包给优化器自动求解的纯技术题，而是需要社会价值共识来决断的伦理选择（究竟更容忍误伤无辜者，还是更容忍漏放高风险者？）。
\end{insightBox}

\subsection{实测验证：完美校准的统一阈值为何制造超过 $5$ 倍的误报率鸿沟？}

我们在配套脚本 \texttt{code/vol7\_experiments.py} 的 \texttt{exp4\_integrated\_\allowbreak gradients\_and\_fairness} 中构造了一个无可挑剔的“理想贝叶斯校准”对照实验：两个群体各有 $N = 10{,}000$ 名个体，群体 A 基础率 $\pi_A = 0.20$，群体 B 基础率 $\pi_B = 0.50$。每名个体的风险分值 $S$ 按 $\operatorname{Beta}(2\pi_g,\, 2(1-\pi_g))$ 采样，真实标签直接按 $Y \sim \operatorname{Bernoulli}(S)$ 生成——换言之，\textbf{分值 $S$ 对两个群体都是数学上 $100\%$ 严格校准的真实后验概率 $\Prob(Y=1\mid S=s, A=g) = s$}！我们对两个群体施以完全相同的统一决策阈值 $\hat{Y} = \mathbb{I}[S \ge 0.50]$，实测统计汇总于表 \ref{tab:fairness-sim}。

\begin{table}[htbp]
  \centering
  \caption{完美校准分值在统一阈值 $t=0.50$ 下的群体公平性实测对比（每组 $N = 10{,}000$）}
  \label{tab:fairness-sim}
  \footnotesize
  \begin{tabular}{@{}lcccccc@{}}
    \toprule
    群体 & 基础率 $\pi_g$ & 精确率 $\mathrm{PPV}_g$ & 误报率 $\mathrm{FPR}_g$ & 真阳性率 $\mathrm{TPR}_g$ & 漏报率 $\mathrm{FNR}_g$ & 恒等式预测 $\mathrm{FPR}$ \\
    \midrule
    群体 A & $0.20$ & $0.703$ & $\mathbf{0.049}$ & $0.454$ & $0.546$ & $0.25\!\times\!0.422\!\times\!0.454 = 0.048$ \\
    群体 B & $0.50$ & $0.752$ & $\mathbf{0.252}$ & $0.755$ & $0.245$ & $1.00\!\times\!0.330\!\times\!0.755 = 0.249$ \\
    \midrule
    \textbf{比值（B/A）} & $2.50\times$ & $1.07\times$ & $\mathbf{5.14\times}$（悬殊!） & $1.66\times$ & $0.45\times$ & 与实测严丝合缝 \\
    \bottomrule
  \end{tabular}
\end{table}

结果触目惊心：即便评分函数完全没有偏见、对所有群体都做到了 $100\%$ 完美概率校准且采用一视同仁的阈值 $t=0.50$，两个群体的 $\mathrm{PPV}$ 相当接近（$0.703$ vs $0.752$），但群体 B 无辜者被错误打上高风险标签的误报率 $\mathrm{FPR}$ 高达 $0.252$，竟是群体 A（$0.049$）的 \textbf{$5.14$ 倍}！与此同时，群体 A 真正需要预警的阳性个体却有高达 $54.6\%$ 被漏报（$\mathrm{TPR}_A = 0.454$ vs $\mathrm{TPR}_B = 0.755$）。把表中的实测数代入定理 \ref{thm:chouldechova} 的恒等式右侧，分别算得 $0.048$ 与 $0.249$，与实测 $\mathrm{FPR}$ 误差仅在有限样本波动之内。如果业务方强行通过为不同群体设置差异化阈值来拉平 $\mathrm{FPR}$ 与 $\mathrm{TPR}$，那么 $\mathrm{PPV}$ 就必然被撕裂（违背校准公平性），甚至在部分法域触及“基于受保护属性的差别待遇”法律争议。

\begin{failBox}{可解释性与公平性实践中的三大常见雷区}
\begin{enumerate}
  \item \textbf{“不知情即公平（Fairness Through Unawareness）”幻觉}：许多工程师以为只要从输入特征表里把“性别”“种族”“年龄”这几列直接 \texttt{drop} 掉，模型就一定公平了。这是极端天真的误解：在现实高维数据中，\textbf{邮政编码、消费时段、设备型号、语言用词习惯}等看似中性的\textbf{代理变量（Proxy Variables）}与受保护属性高度相关，强大非线性模型会毫不费力地隐式重建出受保护属性，而且删掉敏感列后你连事后审计群体误报率差异的抓手都失去了！
  \item \textbf{视觉显著性图的“樱桃采摘”与健全性检验失败}：Adebayo 等人（NeurIPS 2018）提出了一项令解释领域冷汗直冒的\textbf{模型参数随机化健全性检验（Sanity Check）}——把训练好的卷积网络权重逐层重置为纯随机噪声，重新运行多种热门显著性图方法，惊人地发现 Guided Backprop 与 Guided Grad-CAM 生成的热力图依然清晰勾勒出原图中的物体轮廓！它们根本不是在解释“模型学到了什么”，而是在做与网络权重无关的“输入图像边缘检测”。任何显著性图方法上线前都必须先通过权重随机化与标签打乱两项健全性检验。
  \item \textbf{把线性探针的高分当成因果机制}：看到探针准确率达到 $95\%$ 就宣称“我们发现了控制模型道德判断的神经元”，却不做控制任务选择性检验，也不做激活修补因果干预（即抹掉该方向看输出是否真变），极易陷入自我感动的伪解释。
\end{enumerate}
\end{failBox}

\begin{extBox}{反事实因果公平性与 Transformer 机械电路实例（IOI / 归纳头）}
\begin{itemize}
  \item \textbf{反事实公平性（Counterfactual Fairness，Kusner et al., NeurIPS 2017）}：群体统计指标只衡量宏观比例，却无法回答单个个体心中的追问：“如果仅仅改变我的受保护属性 $A$（而保持与 $A$ 无关的个人潜变量 $U$ 不变），模型给我的决策会改变吗？”在结构因果模型（SCM）中，若对任意个体观测上下文 $(x, a)$ 与反事实属性 $a'$ 均满足
  \[
    \Prob\!\big(\hat{Y}_{A \leftarrow a}(U) = y \mid X=x, A=a\big) \;=\; \Prob\!\big(\hat{Y}_{A \leftarrow a'}(U) = y \mid X=x, A=a\big),
  \]
  则称决策 $\hat{Y}$ 具有反事实公平性。其充分条件是 $\hat{Y}$ 仅依赖因果图上不是 $A$ 后代的变量；难点在于现实中很难毫无争议地画出真实的社会因果有向无环图。
  \item \textbf{Transformer 内部真实发现的经典电路}：机械可解释性已经在真实语言模型中完整逆向工程出多条令人惊叹的算法电路：（1）\textbf{归纳头（Induction Heads，Olsson et al., 2022）}：由前一层的“前驱词头（Previous Token Head）”把位置 $i-1$ 的信息写入位置 $i$ 的残差流，再由后一层的归纳头通过 Q--K 匹配寻找历史中相同前缀 $[A][B]\dots[A]$ 并复制 $[B]$，被证实是驱动大模型上下文学习（In-Context Learning）相变的核心双层电路；（2）\textbf{间接宾语识别电路（IOI Circuit，Wang et al., ICLR 2023）}：在 GPT-2 Small 中处理 ``When Mary and John went to the store, John gave a drink to \_\_\_'' 时，研究者用路径修补完整绘出了由 $26$ 个注意力头协同构成的电路——包含\emph{重复词元头}识别出 ``John'' 出现两次、\emph{抑制头（S-Inhibition Heads）}压低 ``John'' 的注意力查询、以及\emph{名称搬运头（Name Mover Heads）}将唯一未被抑制的 ``Mary'' 搬运到输出 logit！
\end{itemize}
\end{extBox}

\section{本章小结与核心公式速查}

本章系统梳理了 AI 可解释性与算法公平性的数学基石。从原始梯度在饱和区的敏感性崩溃出发，我们用多元微积分基本定理严格证明了积分梯度的完备性公理，又通过合作博弈论推导出唯一的 Shapley 归因公式；随后深入网络内部，借助激活修补与稀疏自编码器从“相关性观测”迈向“因果电路解耦”；最后以四行贝叶斯代数证明了 Chouldechova 公平性不可能定理，看清了校准公平与均等赔率在基础率差异下的不可调和张力。表 \ref{tab:ch4-summary} 汇总了本章核心公式与审计要点。

\begin{table}[htbp]
  \centering
  \caption{第 4 章可解释性、机械电路与算法公平性核心公式速查}
  \label{tab:ch4-summary}
  \footnotesize
  \begin{tabularx}{\linewidth}{@{}p{33mm}p{62mm}X@{}}
    \toprule
    \textbf{方法 / 定理} & \textbf{核心公式} & \textbf{核心性质与审计要点} \\
    \midrule
    积分梯度（IG） & $\IG_i(x;x') = (x_i-x'_i)\int_0^1 \partial_i F(x'+\alpha(x-x'))\,d\alpha$ & 满足敏感性、实现不变性与完备性 $\sum \IG_i = F(x)-F(x')$ \\
    中点黎曼求积 & $\alpha_k = (k-0.5)/m$，误差 $O(1/m^2)$ & 可用完备性残差在线自检；警惕单一基线盲区 \\
    Shapley 值 & $\phi_i(v) = \sum_{S\subseteq N\setminus\{i\}} \frac{|S|!(d-|S|-1)!}{d!}[v(S\cup\{i\})-v(S)]$ & 唯一满足效率、对称、哑元、可加；树模型用 TreeSHAP \\
    激活修补效应 & $\mathrm{IE}(\ell,t) = \frac{L_{\text{patch}}(\ell,t)-L_{\text{corr}}}{L_{\text{clean}}-L_{\text{corr}}}$ & 基于 $\mathrm{do}$-干预识别因果电路，克服探针相关陷阱 \\
    稀疏自编码器（SAE） & $\min \E[\|h - W_{\text{dec}}z - b\|_2^2 + \lambda\|z\|_1]$，$\|W_{\text{dec},:,j}\|_2=1$ & 解耦多义叠加激活；必须约束解码器列范数为 $1$ \\
    公平性不可能恒等式 & $\mathrm{FPR}_g = \frac{\pi_g}{1-\pi_g}\cdot\frac{1-\mathrm{PPV}_g}{\mathrm{PPV}_g}\cdot(1-\mathrm{FNR}_g)$ & $\pi_0 \neq \pi_1$ 且非完美时，组内校准与均等赔率不可兼得 \\
    \bottomrule
  \end{tabularx}
\end{table}

\section{分层习题}

\exTier{A 层}{概念核查与基础演练}

\begin{enumerate}[label=\textbf{7.4.\arabic*}]
  \item \textbf{（二元多项式的积分梯度精确手算）} 设双变量函数 $F(x_1, x_2) = x_1^2 x_2 + 3 x_2$，取基线 $x' = (0, 0)^\top$，输入 $x = (2, 1)^\top$。
  \begin{enumerate}
    \item 分别计算端点 $x$ 处的 \texttt{Input} $\times$ \texttt{Gradient} 向量，检查其分量之和是否等于 $F(x) - F(x')$；
    \item 解析计算积分梯度 $\IG_1(x; x')$ 与 $\IG_2(x; x')$，并验算完备性公理 $\IG_1 + \IG_2 = F(x) - F(x')$。
  \end{enumerate}
  \exHint{沿路径 $\gamma(\alpha) = (2\alpha, \alpha)$ 有 $\partial_1 F(\gamma(\alpha)) = 4\alpha^2$，$\partial_2 F(\gamma(\alpha)) = 4\alpha^2 + 3$；直接积分得 $\IG_1 = 8/3$，$\IG_2 = 13/3$，和为 $7 = F(2,1)-F(0,0)$！而 \texttt{Input}$\times$\texttt{Gradient} 之和高达 $15 \neq 7$。}

  \item \textbf{（Chouldechova 恒等式的数值推演）} 某信贷违约预测系统服务两个地域群体，群体 A 的真实违约基础率 $\pi_A = 0.10$，群体 B 的真实违约基础率 $\pi_B = 0.30$。监管要求模型在两组上均保持精确率 $\mathrm{PPV}_A = \mathrm{PPV}_B = 0.60$，且业务要求两组均达到相同的风险捕捉率（$\mathrm{TPR}_A = \mathrm{TPR}_B = 0.80$）。
  \begin{enumerate}
    \item 利用定理 \ref{thm:chouldechova} 的精确恒等式，分别计算群体 A 与群体 B 必然产生的误报率 $\mathrm{FPR}_A$ 与 $\mathrm{FPR}_B$；
    \item 群体 B 信用良好者被误拒的概率是群体 A 的多少倍？如果强行下调群体 B 的阈值以把 $\mathrm{FPR}_B$ 压低到与 $\mathrm{FPR}_A$ 相等，会发生什么数学后果？
  \end{enumerate}
  \exHint{直接代入 $\mathrm{FPR}_g = \frac{\pi_g}{1-\pi_g}\cdot\frac{0.40}{0.60}\cdot 0.80$：群体 A 得 $\frac{1}{9}\times\frac{8}{15} \approx 5.93\%$，群体 B 得 $\frac{3}{7}\times\frac{8}{15} \approx 22.86\%$（近 $3.86$ 倍！）。}
\end{enumerate}

\exTier{B 层}{数学推导与定理证明}

\begin{enumerate}[label=\textbf{7.4.\arabic*}, start=3]
  \item \textbf{（路径归因的唯一性：为何积分梯度选择直线路径？）} 给定任意一条从 $x'$ 到 $x$ 的分段光滑路径 $\gamma: [0,1] \to \R^d$（$\gamma(0)=x', \gamma(1)=x$），都可以定义\textbf{路径积分梯度}
  \[
    \mathrm{PathIG}_i^\gamma(x; x') = \int_0^1 \frac{\partial F(\gamma(\alpha))}{\partial x_i}\,\gamma'_i(\alpha)\,d\alpha.
  \]
  \begin{enumerate}
    \item 证明对\textbf{任意}光滑路径 $\gamma$，完备性公理 $\sum_{i=1}^d \mathrm{PathIG}_i^\gamma(x; x') = F(x) - F(x')$ 依然成立；
    \item 取对称函数 $F(x_1, x_2) = x_1 x_2$、基线 $(0,0)$ 与输入 $(1,1)$，构造一条先走 $x_1$ 轴再走 $x_2$ 轴的折线路径 $\gamma_{\text{L}}$，验证它虽然满足完备性，却彻底破坏了\textbf{对称保持性公理}（请据此解释为何只有直线路径能唯一保持对称性）。
  \end{enumerate}
  \exHint{沿折线 $(0,0)\to(1,0)\to(1,1)$，第一段 $x_2=0$ 故 $\partial_1 F = x_2 = 0 \implies \mathrm{PathIG}_1 = 0$，第二段 $\partial_2 F = 1 \implies \mathrm{PathIG}_2 = 1$，完全把功劳给了后走的特征！}

  \item \textbf{（Shapley 值与积分梯度的退化等价性及差异分析）} 考虑线性函数 $F(x) = w^\top x + b$ 与含二阶交互的二次函数 $G(x) = x_1 x_2$，假设各输入特征相互独立且基线取均值 $x' = \E[X] = \mathbf{0}$。
  \begin{enumerate}
    \item 证明对于任意线性函数 $F(x) = w^\top x + b$，积分梯度 $\IG_i(x; \mathbf{0})$、特征条件期望 Shapley 值 $\phi_i(x)$ 以及 \texttt{Input} $\times$ \texttt{Gradient} 三者\textbf{完全重合}，均等于 $w_i x_i$；
    \item 对二次交互函数 $G(x) = x_1 x_2$（设 $X_1, X_2$ 独立且零均值），分别计算 $\IG_1(x; \mathbf{0})$ 与 Shapley 值 $\phi_1(x)$，比较两者是否仍然相等，并检验三次交互项 $H(x) = x_1^2 x_2$ 上两者是否开始分道扬镳！
  \end{enumerate}
  \exHint{在 $G(x)=x_1 x_2$ 上两者恰好均为 $\frac{1}{2}x_1 x_2$；而在非对称高次项 $H(x)=x_1^2 x_2$ 上，Shapley 值把协同平分（各得 $\frac{1}{2}x_1^2 x_2$），而积分梯度按幂次权重给出 $\IG_1 = \frac{2}{3}x_1^2 x_2$、$\IG_2 = \frac{1}{3}x_1^2 x_2$！}
\end{enumerate}

\exTier{C 层}{算法分析与编程实验}

\begin{enumerate}[label=\textbf{7.4.\arabic*}, start=5]
  \item \textbf{（实证对比中点法 vs 左端点黎曼和的积分梯度收敛阶）} 修改 \texttt{code/vol7\_experiments.py} 中的 \texttt{exp4}：
  \begin{enumerate}
    \item 额外实现\textbf{左端点法}（$\alpha_k = (k-1)/m$）与\textbf{梯形法}，在步数 $m \in \{2, 5, 10, 20, 50, 100, 200\}$ 上记录三者的完备性误差，并在双对数坐标系中拟合各自斜率，验证左端点为 $O(1/m)$ 而中点与梯形均为 $O(1/m^2)$；
    \item 将网络激活函数从光滑的 \texttt{nn.Tanh()} 替换为分段线性的 \texttt{nn.ReLU()}，观察完备性误差收敛阶是否因折点处二阶导数不存在而出现局部抖动。
  \end{enumerate}
  \exHint{光滑函数上中点与梯形法均为严格的斜率 $-2$；遇到 ReLU 折点跨越子区间时，由于泰勒余项在不可微点失效，局部收敛阶会轻微退化。}

  \item \textbf{（在玩具叠加模型上训练稀疏自编码器解耦特征）} 构造合成叠加数据：真实稀疏特征 $c \in \R_{\ge 0}^{16}$ 每个样本仅随机激活 $2$ 个非零分量（$c_j \sim \operatorname{Unif}[0.5, 1.5]$），经随机单位列矩阵 $W_{\text{true}} \in \R^{6 \times 16}$ 压缩投影到 $6$ 维隐藏激活 $h = W_{\text{true}} c \in \R^6$。
  \begin{enumerate}
    \item 按式 \eqref{eq:sae}--\eqref{eq:sae-loss} 用 PyTorch 实现过完备 SAE（$6 \to 24 \to 6$，每步梯度更新后执行 \texttt{W\_dec.data = F.normalize(W\_dec.data, dim=0)}）；
    \item 扫描稀疏系数 $\lambda \in \{0.0, 0.05, 0.15, 0.5\}$，统计每个真实方向与学到的字典列之间的最大余弦相似度均值（Mean Max Cosine Similarity, MMCS），并演示撤掉单位列范数约束后训练如何崩溃。
  \end{enumerate}
  \exHint{合理的 $\lambda \approx 0.05\text{--}0.15$ 可使 MMCS 超过 $0.95$，几乎无损恢复被压缩进 $6$ 维空间里的 $16$ 个原始概念方向！}
\end{enumerate}

\exTier{D 层}{开放探究与科研实战}

\begin{enumerate}[label=\textbf{7.4.\arabic*}, start=7]
  \item \textbf{（高风险医疗分诊系统的公平性与可解释性审计方案）} 某三甲医院计划部署一款“急诊脓毒症（Sepsis）早期预警 AI 系统”，实测发现老年群体基础发病率（$\pi_{\text{old}} = 0.25$）显著高于青年群体（$\pi_{\text{young}} = 0.05$）。
  \begin{enumerate}
    \item 结合定理 \ref{thm:chouldechova} 与临床医学风险不对称性（漏诊脓毒症可能致命，而误报仅导致一次额外血检），论证在此场景下应当优先保障\textbf{组内校准}、\textbf{均等漏报率（Equal Opportunity）}还是\textbf{统计均等}，并说明为何必须放弃另外两项；
    \item 医生抱怨热力图归因难以指导处置，请设计一套面向临床医生的“归因 + 反事实建议 + 健全性审计”三合一解释界面规范。
  \end{enumerate}
  \exHint{脓毒症漏诊代价远高于误报代价，优先确保高灵敏度（低 $\mathrm{FNR}$）或组内校准配合按临床效用定制的阈值；解释界面应给出可验证的生理指标归因（如乳酸、血压趋势）而非黑箱像素块。}

  \item \textbf{（设计一套针对大语言模型“拒绝安全边界”的激活修补实验）} 假设你想定位某开源 7B 对话模型中究竟是哪几层注意力头或 MLP 在识别“有害请求”并触发拒绝回答（``抱歉，我无法协助……''）。请写出完整实验设计方案：
  \begin{enumerate}
    \item 如何构造语义高度匹配但安全属性恰好相反的最小对比对 $(x_{\text{harmful}},\, x_{\text{benign}})$（例如把危险动词替换为良性动词），以防止位置偏移与词长噪声污染修补信号？
    \item 定义定量的拒绝方向打分指标（如拒绝首词元与顺从首词元的 logit 差），说明如何区分“\textbf{检测有害意图}的中层特征”与“\textbf{执行礼貌拒绝模板}的高层输出头”，以及如何用方向消融（Directional Ablation）做双向因果充分性与必要性验证。
  \end{enumerate}
  \exHint{Arditi 等人（NeurIPS 2024）的经典研究发现，许多主流开源 LLM 的拒绝行为竟由残差流中一个近似一维的“拒绝方向”介导——正向注入该方向会让模型拒绝良性问题，正交投影擦除该方向则会瓦解拒绝防线（同时也揭示了单靠浅层对齐的安全脆弱性！）。}
\end{enumerate}
